% Circuito de sinal da Questão 2 em altas frequências: transistor PNP em base comum,
% no modelo T (região sombreada), com C_pi entre emissor e base e C_mu entre coletor e base.
% Compilar com pdflatex; o PDF resultante substitui figuras/fig_q2_hf.pdf.
\documentclass[11pt,border=4pt]{standalone}
% Estilo comum dos esquemáticos do relatório (circuitikz).
% Uso: \documentclass[11pt,border=4pt]{standalone} e, em seguida, % Estilo comum dos esquemáticos do relatório (circuitikz).
% Uso: \documentclass[11pt,border=4pt]{standalone} e, em seguida, % Estilo comum dos esquemáticos do relatório (circuitikz).
% Uso: \documentclass[11pt,border=4pt]{standalone} e, em seguida, \input{estilo_circuitos}.
\usepackage[T1]{fontenc}
\usepackage[utf8]{inputenc}
\usepackage[brazilian]{babel}
\usepackage{lmodern}
\usepackage{amsmath}
\usepackage[american]{circuitikz}
\usetikzlibrary{calc,arrows.meta}

% Cores das figuras originais
\definecolor{anot}{RGB}{0,135,200}      % anotações: valores calculados, setas, resistências vistas
\definecolor{sombra}{RGB}{220,242,253}  % região do transistor nos modelos T e pi
\definecolor{term}{RGB}{62,124,148}     % nomes de terminais destacados (B, C, E)

% Fios com 0,4 pt e corpo dos componentes com 0,8 pt, como nas figuras originais
\ctikzset{bipoles/thickness=2}
.
\usepackage[T1]{fontenc}
\usepackage[utf8]{inputenc}
\usepackage[brazilian]{babel}
\usepackage{lmodern}
\usepackage{amsmath}
\usepackage[american]{circuitikz}
\usetikzlibrary{calc,arrows.meta}

% Cores das figuras originais
\definecolor{anot}{RGB}{0,135,200}      % anotações: valores calculados, setas, resistências vistas
\definecolor{sombra}{RGB}{220,242,253}  % região do transistor nos modelos T e pi
\definecolor{term}{RGB}{62,124,148}     % nomes de terminais destacados (B, C, E)

% Fios com 0,4 pt e corpo dos componentes com 0,8 pt, como nas figuras originais
\ctikzset{bipoles/thickness=2}
.
\usepackage[T1]{fontenc}
\usepackage[utf8]{inputenc}
\usepackage[brazilian]{babel}
\usepackage{lmodern}
\usepackage{amsmath}
\usepackage[american]{circuitikz}
\usetikzlibrary{calc,arrows.meta}

% Cores das figuras originais
\definecolor{anot}{RGB}{0,135,200}      % anotações: valores calculados, setas, resistências vistas
\definecolor{sombra}{RGB}{220,242,253}  % região do transistor nos modelos T e pi
\definecolor{term}{RGB}{62,124,148}     % nomes de terminais destacados (B, C, E)

% Fios com 0,4 pt e corpo dos componentes com 0,8 pt, como nas figuras originais
\ctikzset{bipoles/thickness=2}

\begin{document}
\begin{circuitikz}
  % ---------------------------------------------------------------- colunas (cm)
  \def\xvi{0}      % fonte v_i
  \def\xRSa{0.6}   % início de R_S
  \def\xRE{2.7}    % nó de R_S, R_E e emissor
  \def\xBg{3.1}    % terra da base
  \def\xBE{3.7}    % terminais B e E
  \def\xS{4.1}     % borda esquerda da região sombreada
  \def\xA{5.6}     % fonte controlada e r_e
  \def\xK{7.0}     % C_mu e C_pi
  \def\xSd{8.35}   % borda direita da região sombreada
  \def\xC{8.85}    % terminal C
  \def\xRC{10.05}  % R_C
  \def\xRL{11.85}  % R_L
  \def\xvo{12.65}  % saída v_o
  % ---------------------------------------------------------------- linhas (cm)
  \def\yE{0}       % emissor
  \def\yB{2.2}     % base
  \def\yN{4.6}     % nó de alpha i_e e C_mu
  \def\yC{5.2}     % coletor
  \def\yG{-2.2}    % terras de v_i e R_E
  \def\yBg{1.7}    % terra da base
  \def\yRg{3.2}    % terras de R_C e R_L
  % meias larguras dos símbolos (circuitikz, bipoles/length = 1,4 cm)
  \def\hR{0.21}    % resistor
  \def\hK{0.42}    % placas do capacitor
  \def\hD{0.49}    % losango da fonte controlada
  \def\hV{0.42}    % círculo da fonte de tensão
  % respiros: cada valor deixa cerca de 6 pt de branco entre a tinta do rótulo e a do desenho
  \def\folga{5pt}     % rótulo ao lado do símbolo
  \def\folgaR{3.8pt}  % R_C e R_L (o itálico e o vértice do zigue-zague já afastam o R)
  \def\folgat{7.6pt}  % do centro do terminal (raio 1,6 pt) à base do nome do terminal
  \def\folgao{7.1pt}  % do centro do terminal de saída ao rótulo v_o
  \def\folgas{5.2pt}  % do topo do zigue-zague de R_S à base do rótulo

  % ---------------------------------------------------------------- região do transistor (modelo T)
  \fill[sombra] (\xS,\yE-0.35) rectangle (\xSd,\yC+0.3);

  % ---------------------------------------------------------------- fios e componentes
  % entrada: v_i, R_S, R_E e emissor
  \draw (\xvi,\yE) to[V] (\xvi,\yG);
  \draw (\xvi,\yE) -- (\xRSa,\yE) to[R] (\xRE,\yE) -- (\xK,\yE);
  \draw (\xRE,\yE) to[R] (\xRE,\yG);
  % r_e e C_pi, entre emissor e base
  \draw (\xA,\yE) to[R] (\xA,\yB);
  \draw (\xK,\yE) to[C] (\xK,\yB);
  % base em terra
  \draw (\xBg,\yBg) -- (\xBg,\yB) -- (\xK,\yB);
  % fonte controlada alpha i_e e C_mu, entre base e coletor
  \draw (\xA,\yB) to[cI] (\xA,\yN) -- (\xA,\yC) -- (\xvo,\yC);
  \draw (\xK,\yB) to[C] (\xK,\yN) -- (\xA,\yN);
  % carga: R_C e R_L
  \draw (\xRC,\yC) to[R] (\xRC,\yRg);
  \draw (\xRL,\yC) to[R] (\xRL,\yRg);

  % ---------------------------------------------------------------- terras
  \draw (\xvi,\yG) node[ground]{};
  \draw (\xRE,\yG) node[ground]{};
  \draw (\xBg,\yBg) node[ground]{};
  \draw (\xRC,\yRg) node[ground]{};
  \draw (\xRL,\yRg) node[ground]{};

  % ---------------------------------------------------------------- junções e terminais
  \node[circ] at (\xRE,\yE) {};
  \node[circ] at (\xA,\yB) {};
  \node[circ] at (\xK,\yB) {};
  \node[circ] at (\xA,\yN) {};
  \node[circ] at (\xRC,\yC) {};
  \node[circ] at (\xRL,\yC) {};
  \node[ocirc] at (\xBE,\yE) {};
  \node[ocirc] at (\xBE,\yB) {};
  \node[ocirc] at (\xC,\yC) {};
  \node[ocirc] at (\xvo,\yC) {};

  % ---------------------------------------------------------------- rótulos
  \node[anchor=east, inner sep=0pt, xshift=-\folga] at ({\xvi-\hV},{(\yE+\yG)/2}) {$v_i$};
  \node[anchor=south, inner sep=0pt, yshift=\folgas] at ({(\xRSa+\xRE)/2},\hR) {$R_S$};
  \node[anchor=east, inner sep=0pt, xshift=-\folga] at ({\xRE-\hR},{(\yE+\yG)/2}) {$R_E$};
  \node[anchor=south, inner sep=0pt, yshift=\folgat, text=term] at (\xBE,\yE) {$E$};
  \node[anchor=south, inner sep=0pt, yshift=\folgat, text=term] at (\xBE,\yB) {$B$};
  \node[anchor=south, inner sep=0pt, yshift=\folgat, text=term] at (\xC,\yC) {$C$};
  \node[anchor=east, inner sep=0pt, xshift=-\folga] at ({\xA-\hR},{(\yE+\yB)/2}) {$r_e$};
  \node[anchor=west, inner sep=0pt, xshift=\folga] at ({\xK+\hK},{(\yE+\yB)/2}) {$C_\pi$};
  \node[anchor=east, inner sep=0pt, xshift=-\folga] at ({\xA-\hD},{(\yB+\yN)/2}) {$\alpha i_e$};
  \node[anchor=west, inner sep=0pt, xshift=\folga] at ({\xK+\hK},{(\yB+\yN)/2}) {$C_\mu$};
  \node[anchor=west, inner sep=0pt, xshift=\folgaR] at ({\xRC+\hR},{(\yC+\yRg)/2}) {$R_C$};
  \node[anchor=west, inner sep=0pt, xshift=\folgaR] at ({\xRL+\hR},{(\yC+\yRg)/2}) {$R_L$};
  \node[anchor=west, inner sep=0pt, xshift=\folgao] at (\xvo,\yC) {$v_o$};
\end{circuitikz}
\end{document}
